This work reproduced the unsupervised edge detection pipeline of Zotin et al. \cite{zotin2018edge} and validated it on 484 multiparametric volumes from the \gls{msd} Brain Tumor task \cite{antonelli2022medical}, a scale two orders of magnitude larger than the 30 hand-picked images used in the original study. This large-scale validation revealed a weakness that a small curated dataset does not show: mean Sensitivity dropped to 10.7\%, mainly because of confusion between tumor and \gls{csf} hyperintensity on T2-weighted images, and because a static central-slice selection did not even intersect the annotated lesion for 8.1\% of patients. The proposed optimized pipeline addressed both causes at the input level, without modifying the underlying \gls{fcm} clustering or edge-extraction logic, by combining an unsupervised dynamic slice search with a multimodal transition to \gls{flair} sequences and a spatially tolerant evaluation metric. This yielded a 132\% increase in mean Sensitivity and a 46\% increase in mean \gls{fom}, at a small additional computational cost. Overall, this suggests that the reliability of an otherwise unchanged unsupervised, intensity-based segmentation method depends as much on the physical contrast mechanism and the spatial cross-section it is given as on the clustering algorithm itself: changing the input modality and the slice-selection criterion recovered most of the performance gap, while the clustering and edge-extraction stages stayed the same. At the same time, the shared zero-\gls{fom} failure on the same patient across both pipelines shows that a global, volume-statistics-based proxy for tumor localization, however effective on average, is not a substitute for an explicit spatial or shape prior, and can still fail outright on small or eccentrically located lesions. Future work should target this remaining gap directly, for instance through deterministic or atlas-based centroid initialization for \gls{fcm} \cite{benson2016brain, szilagyi2003mri}, morphological skull-stripping refinements to further constrain the clustering search space \cite{hooda2014brain}, a purely spatial criterion that discounts centrally-located clusters under the prior that the lateral ventricles, unlike most tumors, sit near the brain's geometric center -- a lower-cost alternative to the \gls{flair} transition when only T2 is available -- or a learned, supervised localization step to replace the unsupervised hyperintensity proxy for slice and region selection. More generally, this study suggests that unsupervised medical segmentation pipelines should be evaluated, whenever feasible, on large and heterogeneous clinical datasets rather than on small curated samples, since the latter can hide failure modes that only show up at scale. More broadly, the entire pipeline studied here belongs to a class of unsupervised, intensity-based methods that has been substantially superseded in the medical imaging literature by supervised deep learning \cite{magadza2021survey}. Architectures such as U-Net \cite{ronneberger2015unet} and its self-configuring descendant nnU-Net \cite{isensee2021nnunet} learn, directly from labeled examples, hierarchical spatial features that jointly encode intensity, shape, and anatomical context -- precisely the spatial and anatomical prior that this study found missing from both the \gls{fcm} cluster-selection heuristic and the volume-statistics-based slice selection. Trained end-to-end on datasets such as the \gls{msd} Brain Tumor task used here \cite{antonelli2022medical}, these models require no hand-tuned brain mask, no fixed number of clusters, and no brightness-based rule to disambiguate the tumor from \gls{csf}: the segmentation boundary is learned jointly with the features that define it. As a concrete reference point, an nnU-Net variant placed first in the BraTS 2020 challenge with Dice scores of 0.890, 0.851, and 0.820 for whole tumor, tumor core, and enhancing tumor respectively \cite{isensee2020brats} -- figures well beyond what intensity-only clustering can reach on the same task. The corresponding cost is that these methods require a labeled training set and non-trivial training-time compute, whereas the pipeline studied here needs no training data at all and remains fully interpretable at every stage -- which keeps it relevant as a fast, equipment-agnostic baseline, an annotation-assisting tool, or a teaching example of classical image processing -- but for raw segmentation and edge-detection accuracy on annotated clinical datasets, deep learning is, at this point, the dominant and recommended approach.
